
Communication Accommodation Theory predicts that linguistic convergence increases engagement. This study tests this prediction in human-AI dialogue through analysis of formality accommodation in 111,039 conversations from the Anthropic hh-rlhf dataset. Using a DeBERTa-based formality ranker, accommodation is measured as the formality delta between human inputs and AI responses, and its relationship to conversation continuation is examined.

Three findings emerge. First, AI systems exhibit formality accommodation (r=0.152, p<0.001, n=111,039), adapting response formality to match human input. Second, this accommodation is asymmetric: AI-to-human adaptation is approximately 11 times stronger than human-to-AI adaptation (r=0.152 vs r=0.013). Third, the accommodation-engagement relationship is non-linear. Conversations with moderate formality accommodation show 71.4\% continuation rates, compared to 65.9\% for high accommodation and 60.9\% for low accommodation.

This inverted-U pattern challenges the assumption that more accommodation yields better outcomes. Excessive convergence may be perceived as artificial mimicry. These correlational findings suggest that conversational AI designers should consider calibrating accommodation rather than maximizing it.

